% Part: first-order-logic
% Chapter: models-theories
% Section: expressing-relations

\documentclass[../../../include/open-logic-section]{subfiles}

\begin{document}

\olfileid{fol}{mat}{exr}

\olsection{Nyatakake Relasi ing \article{structure}
  \printtoken{S}{structure}}

\begin{explain}
Salah siji piguna pokok !!{formula}s yaiku nyatakake sipat
lan relasi ing !!a{structure}~$\Struct M$ lumantar konsep-konsep
dhasar ing basa~$\Lang L$ saka~$\Struct M$. Tegese mangkene:
!!{domain} saka $\Struct M$ iku himpunan objek.
!!{constant}s, !!{function}s, lan !!{predicate}s
diinterpretasi ing~$\Struct M$ minangka objek-objek
ing~$\Domain M$, fungsi-fungsi ing~$\Domain M$,
lan relasi-relasi ing~$\Domain M$. Upamane, yen
$\Obj A^2_0$ ana ing $\Lang L$, mula $\Struct M$
menehi interpretasi relasi~$R = \Assign{{\Obj
  A^2_0}}{M}$ marang simbol kasebut. Banjur formula
$\Atom{\Obj A^2_0}{\Obj v_1, \Obj v_2}$
\emph{nyatakake} relasi iku dhewe, kanthi teges:
yen pangaji variabel~$s$ masangake $\Obj v_1$
karo $a \in \Domain{M}$ lan $\Obj v_2$
karo $b \in \Domain M$, mula
\[
Rab \text{\quad yen lan mung yen\quad} \Sat{M}{\Atom{\Obj A^2_0}{\Obj v_1, \Obj v_2}}[s].
\]
Gatekna yen ing kene kita kudu nglebokake pangaji
variabel: kita ora bisa mung kandha
"$Rab$ yen lan mung yen $\Sat{M}{\Atom{\Obj A^2_0}{a, b}}$"
amarga $a$ lan $b$ dudu simbol-simbol basa kita:
loro-lorone iku !!{element}s saka~$\Domain{M}$.

Amarga kita ora mung nduweni !!{formula}s atomik,
nanging uga bisa nggabungake nganggo konektif
logis lan kuantor, !!{formula}s kang luwih rumit
bisa nemtokake relasi liyane kang ora langsung
kalebu ing~$\Struct M$. Kita kepengin ngerti
carane nindakake iku, lan mligi, relasi endi
kang bisa kita temtokake ing !!a{structure}.
\end{explain}

\begin{defn}
Upamakna $!A(\Obj v_1,\dots, \Obj v_n)$ iku
!!a{formula} saka $\Lang L$ kang mung $\Obj v_1$,
\dots, $\Obj v_n$ bisa katon bebas, lan
$\Struct M$ iku !!a{structure} kanggo~$\Lang L$.
$!A(\Obj v_1,\dots, \Obj v_n)$
\emph{nyatakake relasi}~$R \subseteq \Domain M^n$
yen lan mung yen
\[
Ra_1\dots a_n \text{\quad yen lan mung yen\quad} \Sat{M}{\Atom{!A}{\Obj
    v_1,\dots, \Obj v_n}}[s]
\]
kanggo sembarang pangaji variabel~$s$ kang
$s(\Obj v_i) = a_i$ ($i = 1, \dots, n$).
\end{defn}

\begin{ex}
Ing model standar aritmetika~$\Struct N$, !!{formula} $\Obj
v_1 < \Obj v_2 \lor \eq[\Obj v_1][\Obj v_2]$
nyatakake relasi $\le$ ing~$\Nat$.
!!{formula} $\eq[\Obj v_2][\Obj v_1']$
nyatakake relasi penerus, yaiku relasi $R \subseteq
\Nat^2$ kang $Rnm$ bener yen $m$ iku penerus saka~$n$.
Formula $\eq[\Obj v_1][\Obj v_2']$
nyatakake relasi pendahulu. !!{formula}s
$\lexists[\Obj v_3][(\eq/[\Obj v_3][\Obj 0] \land
  \eq[\Obj v_2][(\Obj v_1 + \Obj v_3)])]$ lan
$\lexists[\Obj
  v_3][\eq[(\Obj v_1 + {\Obj v_3}')][v_2]]$
loro-lorone nyatakake relasi $\Obj <$.
Iki ateges simbol predikat~$<$ sejatine
keluwihan ing basa aritmetika; relasi iku
bisa ditemtokake.
\end{ex}

\begin{explain}
Gagasan iki ora mung narik kawigaten kanggo
!!{structure}s tartamtu, nanging uga lumaku
umume nalika kita nggunakake basa kanggo
njlentrehake model utawa model-model kang
dikarepake, yaiku nalika kita ngrembug teori.
Teori-teori iki kerep mung ngemot sawetara
!!{predicate}s minangka simbol dhasar, nanging
ing ranah kang dijlentrehake kerep ana akeh
relasi liyane kang perane wigati. Yen relasi-relasi
liyane iku bisa dinyatakake kanthi sistematis
lumantar relasi kang dadi interpretasi
!!{predicate}s dhasar ing basa kasebut,
kita kandha yen relasi-relasi iku bisa
\emph{ditemtokake} ing basa kasebut.
\end{explain}

\begin{prob}
Temokna !!{formula}s ing $\Lang L_A$ kang
nemtokake relasi-relasi ing ngisor iki:
\begin{enumerate}
\item $n$ dumunung ing antarane $i$ lan $j$;
\item $n$ mbagi $m$ tanpa turahan (yaiku, $m$
  iku kelipatan saka $n$);
\item $n$ iku wilangan prima (yaiku, luwih gedhe tinimbang siji lan ora ana
  wilangan saliyane $1$ lan $n$ kang mbagi~$n$
  tanpa turahan). Cathetan koreksi sumber SOL6-F009: syarat luwih gedhe tinimbang siji ditambah amarga parenthesis sumber uga dituruti dening wilangan siji.
\end{enumerate}
\end{prob}

\begin{prob}
Upamakna formula $!A(\Obj v_1, \Obj v_2)$
nyatakake relasi $R
\subseteq \Domain M^2$ ing
!!a{structure}~$\Struct M$. Temokna formula-formula
kang nyatakake relasi-relasi ing ngisor iki:
\begin{enumerate}
\item invers $R^{-1}$ saka $R$;
\item produk relatif $R \mid R$;
\end{enumerate}
Apa bisa nemokake cara kanggo nyatakake $R^+$,
tutupan transitif saka~$R$?
\end{prob}

\begin{prob}
Upamakna $\Lang{L}$ iku basa kang mung ngemot
simbol predikat loro panggonan $<$ (tanpa
!!{constant}s, !!{function}s, utawa !!{predicate}s
liyane---kajaba mesthi~$\eq$). Upamakna
$\Struct{N}$ iku struktur kang
$\Domain{N} = \Nat$, lan
$\Assign{<}{N} = \Setabs{\tuple{n,m}}{n <
  m}$. Buktekna ing ngisor iki:
\begin{enumerate}
\item $\{ 0 \}$ bisa ditemtokake ing $\Struct{N}$;
\item $\{ 1 \}$ bisa ditemtokake ing $\Struct{N}$;
\item $\{ 2 \}$ bisa ditemtokake ing $\Struct{N}$;
\item kanggo saben $n \in \Nat$, himpunan
  $\{ n \}$ bisa ditemtokake ing $\Struct{N}$;
\item saben subhimpunan winates saka $\Domain{N}$
  bisa ditemtokake ing $\Struct{N}$;
\item saben subhimpunan kofinit saka $\Domain{N}$
  bisa ditemtokake ing $\Struct{N}$
  (ing ngendi $X \subseteq \Nat$ iku kofinit
  yen lan mung yen $\Nat \setminus X$ winates).
\end{enumerate}
\end{prob}


\end{document}
